% Partie : théorie de la démonstration
% Chapitre : normalisation
% Section : introduction

\documentclass[../../../include/open-logic-section]{subfiles}

\begin{document}

\olfileid[fr]{pt}{nor}{int}

\olsection{Introduction}

Si vous avez démontré l'implication~$!A \lif !B$,
vous pouvez l'employer dans une !!{proof} de~$!B$
au moyen de~\Elim{\lif}, pourvu que vous disposiez
aussi d'une !!{proof}~$\delta_1$ de~$!A$.
La preuve de~$!A \lif !B$ peut se terminer par
\Intro{\lif}, qui transforme une !!{proof}~$\delta_2$
de~$!B$ sous l'hypothèse ouverte~$!A$ en une
preuve de~$!A \lif !B$. Dans ce cas, la
!!{proof} finale a la forme
\[
\AxiomC{$\Discharge{!A}{x}$}
\RightLabel{$\delta_2$}
\DeduceC{$!B$}
\DischargeRule{\Intro{\lif}}{x}
\UnaryInfC{$!A \lif !B$}
\AxiomC{}
\RightLabel{$\delta_1$}
\DeduceC{$!A$}
\RightLabel{\Elim{\lif}}
\BinaryInfC{$!B$}
\DisplayProof
\]
Réemployer une preuve déjà obtenue de~$!A \lif !B$
peut faciliter votre travail ; l'arbre obtenu comporte
néanmoins un détour. On introduit $!A \lif !B$ pour
l'éliminer aussitôt. Il devrait être possible de
démontrer~$!B$ plus directement.

C'est effectivement possible. Les détours de ce
genre --- une règle d'introduction suivie de la
règle d'élimination correspondante --- peuvent être
supprimés en déduction naturelle. Cette transformation
est appelée \emph{normalisation} ; elle donne une
!!{proof} \emph{normale}, ou \emph{en forme normale}.
Des éliminations de disjonction ou d'existentielle
peuvent séparer les deux règles d'un détour : les
sections suivantes étendent la définition pour
traiter également ces cas.

Comme l'élimination des coupures en calcul des
séquents, la démonstration de ce résultat exige
de nombreux cas. Le traitement classique de
\Log{N1c} est plus délicat que le traitement
intuitionniste de \Log{N1i}, qui sera développé
dans ce chapitre. De même que l'élimination des
coupures montre que la règle \Cut{} n'ajoute
aucun séquent démontrable au calcul, la
normalisation montre que les détours n'ajoutent
aucune conclusion démontrable.

Les deux résultats présentent d'autres analogies.
Une preuve normale peut être beaucoup plus longue
qu'une preuve plus courte de la même conclusion
comportant des détours. De même, une preuve avec
coupures peut être beaucoup plus courte que toute
preuve sans coupure du même séquent.
La normalisation intuitionniste conduit aussi à
la propriété de la sous-formule, comme l'élimination
des coupures en calcul des séquents. Cette propriété
se rapporte aux formules de la conclusion
\emph{et des hypothèses ouvertes}. Au premier
ordre, elle tient compte des instances des
sous-formules quantifiées. L'absurde peut faire
exception lorsque la négation est primitive.












\end{document}
